\documentclass[titlepage]{ctexart}

\usepackage{amsmath}
\usepackage{amssymb}
% \usepackage{booktabs}
\usepackage{caption}

\title{第六章~概率}
\author{dovego}
\date{\today}

\begin{document}

\maketitle


\section{随机事件的条件概率}

\subsection{条件概率的概念}

我们称
\[
P(B|A) = \frac{P(AB)}{P(A)}
\]
为在事件 $A$ 发生的条件下事件 $B$ 发生的\textbf{条件概率}.

可以通过Venn图来帮助理解.

\vspace{1cm}

\noindent \textbf{思考交流}

1. 概率 $P(B|A)$ 与 $P(AB)$ 有什么区别和联系?

答: (1) $P(B|A) \geqslant P(AB)$. 因为 $P(A) \in (0,1]$.

(2) 求 $P(B|A)$ 时考虑的样本空间中的样本点为事件 $A$ 发生时的样本点, 而求 $P(AB)$ 时考虑的样本空间为整个样本空间 $\varOmega$. 前者小于后者.

2. 若 $B$ 和 $C$ 是两个互斥事件, 则 $P[(B \cup C)|A] = P(B|A) + P(C|A)$是否成立?

答: 成立. 因为 $B$ 与 $C$ 互斥, 所以 $AB$ 与 $AC$ 互斥. 由互斥事件的加法原理
\begin{align*}
    P(B|A) + P(C|A) &= \frac{P(AB)}{P(A)} + \frac{P(AC)}{P(A)} \\
                    &= \frac{P(AB)+P(AC)}{P(A)} \\
                    &= \frac{P[A(B \cup C)]}{P(A)} \\
                    &= P[(B \cup C)|A].
\end{align*}
该公式还可以推广到 $n$ 个事件的情况.

\subsection{乘法公式与事件的独立性}

乘法公式:
\[
P(AB) = P(A)P(B|A).
\]
利用它可以计算两个事件同时发生的概率. 

\textbf{拓展} \quad
(乘法公式的推广)若 $P(A_{1}A_{2}\cdots A_{n}) > 0$, 则
\[
P(A_{1}A_{2}\cdots A_{n}) = P(A_{1})P(A_{2}|A_{1})P(A_{3}|A_{1}A_{2})\cdots P(A_{n}|A_{1}A_{2}\cdots A_{n-1}).
\]

下面的\textbf{例 3} 就对该公式在 $n = 3$ 时进行了推导.

\vspace{1cm}

\textbf{例 3} \quad 已知口袋中有 3 个黑球和 7 个白球, 这 10 个球除颜色外完全相同.

(2) 从中不放回地摸球, 每次各摸一球, 求第三次才摸到黑球的概率. 

\textbf{解} \quad 设事件 $A_{i}$ 表示``第 $i$ 次摸到黑球'', 事件 $A$ 表示``第三次才摸到黑球'',则 $A = \overline{A_{1}A_{2}}A_{3}$.

由题意知 $P(\overline{A_{1}}) = \dfrac{7}{10}, P(\overline{A_{2}}|\overline{A_{1}}) = \dfrac{6}{9}, P(A_{3}|\overline{A_{1}A_{2}}) = \dfrac{3}{8}$.

于是, 根据乘法公式有
\begin{align*}
    P(\overline{A_{1}A_{2}}A_{3}) &= P(\overline{A_{1}A_{2}})P(A_{3}|\overline{A_{1}A_{2}}) \\
                                  &= P(\overline{A_{1}})P(\overline{A_{2}}|\overline{A_{1}})P(A_{3}|\overline{A_{1}A_{2}}) \\
                                  &= \frac{7}{10} \times \frac{6}{9} \times \frac{3}{8} \\
                                  &= \frac{7}{40}.
\end{align*}

所以不放回地摸球, 第三次才摸到黑球的概率为 $\dfrac{7}{40}$.

\vspace{1cm}

若事件 $A$ 与事件 $B$ 相互独立, 则应有 $P(B) = P(B|A)$. 由此可以证得
\[
\text{事件} A \text{与事件} B \text{相互独立} \Longleftrightarrow P(AB) = P(A)P(B).
\]

但当三个事件 $A,B,C$ \textbf{两两独立}时, 等式 $P(ABC) = P(A)P(B)P(C)$ 一般不成立!!!只有说 $A,B,C$ 三个事件\textbf{相互独立} 时才满足该式
(三者相互独立 $\Rightarrow$ 两两独立, 反之不一定成立!!!).

要计算三个事件同时发生的概率, 一般需要利用条件概率计算, 例如
\[
P(ABC) = P(A)P(B|A)P(C|AB).
\]

\textbf{例 4} 告诉我们, 要判断事件 $A, B$ 是否独立, 有以下两种方法:

(1) 判断 $P(B|A) = P(B)$ 是否成立;

(2) 判断 $P(AB) = P(A)P(B)$ 是否成立.

\textbf{补充} \quad
常见的独立事件: 有放回的摸球试验, 掷骰子, 掷硬币, 产品质量检测.

\subsection{全概率公式}

设 $\varOmega$ 是试验 $E$ 的样本空间,$B_{1},B_{2},\dots,B_{n}$ 为样本空间 $\varOmega$ 的一组事件, 若

(1) $B_{i}B_{j} = \varnothing$, 其中 $i \neq j (i,j=1,2,\dots,n)$,

(2) $B_{1}\cup B_{2}\cup \cdots \cup B_{n} = \varOmega$,

则称 $B_{1},B_{2},\dots,B_{n}$ 为样本空间 $\varOmega$ 的一个\textbf{划分}.

\textbf{全概率公式} \quad 设 $B_{1},B_{2},\dots,B_{n}$ 为样本空间 $\varOmega$ 的一个划分, 若 $P(B_{i}) > 0(i=1,2,\dots,n)$, 则对任意一个事件 $A$ 有
\[
P(A) = \sum_{i=1}^n P(B_{i})P(A|B_{i}).
\]
称上式为全概率公式.

这个公式告诉我们, 事件 $A$ 发生的概率是事件 $A$ 在这些``原因'' $(B_{i})$ 下发生的条件概率的平均, 可以看成``由原因推结果''.

\textbf{贝叶斯公式} \quad 设$B_{1},B_{2},\dots,B_{n}$ 为样本空间 $\varOmega$ 的一个划分, 若 $P(A)>0,P(B_{i})>0(i=1,2,\dots,n)$,则
\[
P(B_{i}|A) = \frac{P(B_{i})P(A|B_{i})}{\sum\limits_{j=1}^{n}P(B_{j})P(A|B_{j})}.
\]

称上式为贝叶斯公式. 

这个公式的思想是``执果溯因'', 已知事件 $A$ 发生的情况下, 寻找每个原因的概率.

全概率公式和贝叶斯公式都是新教材添加的内容, 说明对概率的理解要更加深入.


\section{离散型随机变量及其分布列}

\subsection{随机变量}

将样本空间中的每一个样本点都和一个确定的数值(实数)对应起来, 在这个对应下, 数值会随着试验结果的变化而变化. 
将这种数值随着试验结果的变化而变化的量称为\textbf{随机变量}, 常用字母 $X, Y, \xi, \eta$ 等来表示.  

随机变量 $X$ 的一个取值代表一个试验结果. 随机变量一些取值构成的集合表示一个随机事件, 例如: $\{X = 6\}$ 表示事件``掷出的点数为6'', $\{X < 3\}$ 表示事件``掷出的点数小于 3 ''.

\subsection{离散型随机变量的分布列}

取值能够一一列举出来(不论是否有限)的随机变量叫做\textbf{离散型随机变量}.
若离散型随机变量 $X$ 的取值为 $x_{1},x_{2},\dots,x_{n},\dots$, 随机变量 $X$ 取 $x_{i}$ 的概率为 $p_{i}(i=1,2,\dots,n,\dots)$, 记作
\begin{equation}
P(X=x_{i}) = p_{i}(i=1,2,\dots,n,\dots).\label{equ:e1}
\end{equation}

\eqref{equ:e1}式也可以列成表, 如表\ref{tab:t1}:

\begin{table}[htbp]
    \centering
    \caption{}\label{tab:t1}
        \begin{tabular}{c|c|c|c|c|c}
            \hline
            $x_{i}$ & $x_{1}$ & $x_{2}$ & $\cdots$ & $x_{n}$ & $\cdots$ \\
            \hline
            $P(X=x_{i})$ & $p_{1}$ & $p_{2}$ & $\cdots$ & $p_{n}$ & $\cdots$ \\
            \hline
        \end{tabular}
\end{table}

表\ref{tab:t1} 或\eqref{equ:e1} 称为\textbf{离散型随机变量 $\boldsymbol{X}$ 的分布列}, 简称\textbf{$\boldsymbol{X}$ 的分布列}.

显然:

(1) $p_{i} > 0(i=1,2,\dots,n,\dots)$;

(2) $p_{1} + p_{2} + \dots + p_{n} + \dots = 1$.

称随机变量 $X$ 满足这一分布列, 记作
\[
X \sim \left(
\begin{array}{c c c c c}
    x_{1} & x_{2} & \cdots & x_{n} & \cdots \\
    p_{1} & p_{2} & \cdots & x_{n} & \cdots
\end{array}
\right).
\]

\textbf{伯努利试验} \quad 每次试验只有两个对立结果, 可以分别称为``成功''和``失败'', 每次``成功''的概率均为 $p$, 每次``失败''的概率均为 $1-p$.

若随机变量 $X$ 的分布列如表\ref{tab:t2}:

\begin{table}[htbp]
    \centering
    \caption{}\label{tab:t2}
    \begin{tabular}{c|c|c}
        \hline
        X & 1 & 0 \\
        \hline
        P & p & q \\
        \hline
    \end{tabular}
\end{table}

其中 $0< p <1, q = 1 - p$, 那么称离散型随机变量 $X$ 服从参数为 $p$ 的\textbf{两点分布}(又称\textbf{0--1 分布}或\textbf{伯努利分布}).
购买彩票是否中奖, 婴儿的性别, 投篮是否命中等, 都可以用两点分布来描述.


\section{离散型随机变量的均值与方差}

\subsection{离散型随机变量的均值}

我们称
\[
EX = x_{1}p_{1} + x_{2}p_{2} + \cdots + x_{i}p_{i} + \cdots + x_{n}p_{n}
\]
为随机变量 $X$ 的\textbf{均值}或\textbf{数学期望}(简称\textbf{期望}). 数学期望 $EX$ 刻画了随机变量取值的``中心位置'', 但其并不能完全决定随机变量的性质.
随机变量的分布完全确定了它的均值, 不同的分布可以有相同的均值.

两点分布的均值为
\begin{align*}
    EX &= 0 \cdot P(X=0) + 1 \cdot P(X=1) \\
       &= 0 \cdot (1-p) + 1 \cdot p \\
       &= p.
\end{align*}

\subsection{离散型随机变量的方差}

我们知道, 样本方差(初中知识, 也可参看必修一第六章\S{4.4.1}样本的数字特征)可以度量一组样本数据的离散程度. 对于随机变量, 也可以借鉴类似的思想.

我们称 $DX$ 为随机变量 $X$ 的\textbf{方差}, 其算术平方根 $\sqrt{DX}$ 为随机事件 $X$ 的\textbf{标准差}, 记作 $\sigma_{X}$.
\begin{align*}
    DX &= E(X - EX)^2 = \sum_{i=1}^{n}(x_{i} - EX)^2p_{i} \\
       &= E(X^2 - 2XEX + (EX)^2) = EX^2 - (EX)^2.
\end{align*}

方差 $DX$ 为这些偏离程度的加权平均. $DX$ 和 $\sigma_{X}$ 都反映了随机变量的取值偏离均值的平均程度. 

计算离散型随机变量 $X$ 的方差 $DX$ 的一般过程如下:

(1) 计算每一个 $P(X=i) = p_{i}$;

(2) 计算期望 $EX = \sum\limits_{i=1}^{n}x_{i}p_{i}$;

(3) 计算 $DX = E(X - EX)^2 = \sum\limits_{i=1}^{n}(x_{i}-EX)^2p_{i}$.

两点分布的方差为：
\[
DX = E(X-EX)^2 = p(1-p).
\]


\section{二项分布与超几何分布}

\subsection{二项分布}

$n$ 重伯努利试验: 在相同的条件下, 进行 $n$ 次伯努利试验, 且每次试验相互独立.

\textbf{二项分布} \quad 一般地, 在 $n$ 重伯努利试验中, 用 $X$ 表示这 $n$ 次试验中成功的次数, $p$ 为每次试验成功的概率, 则 $X$ 的分布列可以表示为
\[
P(X = k) = C_{n}^k p^k (1-p)^{n-k}(k = 0,1,2,\dots,n).
\]

若随机变量 $X$ 的分布列如上所述, 则称 $X$ 服从参数为 $n,p$ 的\textbf{二项分布}, 简记为 $X\sim B(n,p)$. 两点分布就是 $n = 1$ 时的二项分布.

二项分布的期望与方差为:
\begin{align*}
    EX &= \sum_{k=0}^{n}kC_{n}^{k}p^k(1-p)^{n-k} = np, \\
    DX &= np(1-p).
\end{align*}

当 $n = 1$ 时就是两点分布的期望与方差.

$EX$ 的证明如下:

由 $kC_{n}^{k} = nC_{n-1}^{k-1}$, 得
\[
EX = n\sum_{k=1}^{n}C_{n-1}^{k-1}p^k(1-p)^{n-k} = np\sum_{k=1}^{n}C_{n-1}^{k-1}p^{k-1}(1-p)^{(n-1)-(k-1)}.
\]

令 $m = k - 1$, 由二项式定理得
\[
EX = np\sum_{m=0}^{n-1}C_{n-1}^{m}p^m(1-p)^{(n-1)-m} = np[p + (1 - p)]^{n-1} = np. 
\]

$DX$ 可由 $DX = EX^2 - (EX)^2$ 证得, 略为复杂.

\subsection{超几何分布}

一般地, 设有 $N$ 件产品, 其中有 $M$ 件次品. 从中任取 $n(n \leqslant N)$ 件产品, 用 $X$ 表示取出的 $n$ 件产品中次品的件数, 那么
\[
P(X = k) = \frac{C_{M}^{k}C_{N-M}^{n-k}}{C_{N}^{n}}, \text{max}\{0, n - (N - M)\} < p < \text{min}\{n, M\}.
\]
其中 $n \leqslant N, M \leqslant N, N \in \mathbf{N}_+$.

若随机变量 $X$ 的分布列如上所述, 则称 $X$ 服从参数为 $N, M, n$ 的\textbf{超几何分布}.

超几何分布的期望 $EX$ 为:
\[
EX = \sum_{k=m}^{r}k\frac{C_{M}^{k}C_{N-M}^{n-k}}{C_N^n} = \frac{nM}{N}.
\]
其中 $m = \text{max}\{0, n - (N - M)\}, n = \text{min}\{n, M\}$. 

$EX$ 的证明如下:

\romannumeral1)先证 $\sum\limits_{k=m}^{r}C_{M-1}^{k-1}C_{N-M}^{n-k} = C_{N-1}^{n-1}$.

右边表示在 $N - 1$ 件产品中任选 $n - 1$ 件的总抽法数. 
这是一个组合问题, 现在我们将这 $N - 1$ 件产品分为 $M - 1$ 件次品和 $N - M$ 件正品. 
那么原先的组合问题就可以分解为 $n-1$ 件中``恰好抽到 $m$ 件次品''``恰好抽到 $m + 1$ 件次品''$\dots$``恰好抽到 $r$ 件次品''这若干个互斥事件.
于是原问题的总抽法数便等于这若干个互斥事件的抽法数之和. 若 $n - 1$ 件中抽到 $k$ 件次品, 则有 $C_{M-1}^{k}C_{N-M}^{n-k}$ 种方法.
因此: 分情况累加的总抽法(左边) = 不分情况直接抽的总抽法(右边). 即
\[
\sum\limits_{k=m}^{r}C_{M-1}^{k-1}C_{N-M}^{n-k} = C_{N-1}^{n-1}.
\]

\romannumeral2)再证期望公式.

当 $m > 0$ 时, 由 $kC_M^k = MC_{M-1}^{k-1}, C_N^n = \dfrac{n}{N}C_{N-1}^{n-1}$ 得
\[
EX = \frac{nM}{N} \cdot \frac{1}{C_N^n} \cdot \sum_{k=m}^{r}C_{M-1}^{k-1}C_{N-M}^{n-k} = \frac{nM}{N}.
\]

当 $m = 0$ 时, 类似可证明结论依然成立. 

\subsection{正态分布}

正态分布是最常见、最重要的连续随机变量的分布, 是刻画误差分布的重要模型, 因此也称为\textbf{误差分布}. 记作 $X \sim N(\mu, \sigma^2)$.
其中 $EX = \mu, DX = \sigma^2$. 其分布密度函数的解析式为:
\[
\varphi_{\mu, \sigma} = \frac{1}{\sqrt{2\pi}\sigma^2}e^{-\frac{(x-\mu)^2}{2\sigma^2}}, x \in (-\infty, +\infty).
\]

函数图像关于 $x = \mu$ 对称, 随着 $\mu$ 的变化而平移. 当 $\mu$ 确定时, $\sigma$ 越小, 总体分布越集中; $\sigma$ 越大, 总体分布越分散.

正态分布随机变量 $X$ 在区间 $(a, b]$ 上取值的概率 $P(a < X \leqslant b)$ 就是分布曲线在区间 $(a, b]$对应的曲边梯形面积的值. 特别地, 

(1) $P(\mu - \sigma < X < \mu + \sigma) = 0.683,$

(2) $P(\mu - 2\sigma < X < \mu + 2\sigma) = 0.954,$

(3) $P(\mu - 3\sigma < X < \mu + 3\sigma) = 0.997.$

通常认为服从正态分布的随机变量 $X$ 只取区间 $(\mu - 3\sigma,\mu + 3\sigma]$ 之间的值, 称为\textbf{3$\sigma$原则}. 
并认为 $X$ 在此区间之外取值几乎不可能发生, 是小概率事件.



\end{document}